% ===========================================================================
% Abstract (250 words max)
% Structure: Problem -> Gap -> Approach -> Results -> Impact
% ===========================================================================

% TODO: Write abstract
%
% Suggested structure:
% [Problem] I/O bottlenecks are a leading cause of performance degradation
% in HPC applications, yet diagnosing root causes and providing actionable
% fixes remains a manual, expert-intensive process.
%
% [Gap] Existing tools either detect bottlenecks without recommendations
% (AIIO, Drishti) or generate recommendations without learned detection
% (IOAgent, ION). No system combines ML-based detection with grounded
% code-level recommendations.
%
% [Approach] We present \system{}, a hybrid ML+LLM system that combines
% trained classifiers for multi-label bottleneck detection with a
% retrieval-augmented LLM for code-level fix recommendations. The ML
% component is trained on benchmark-verified gold labels and Drishti-derived
% silver labels from 1.37M Darshan logs collected on ALCF Polaris. SHAP
% attributions identify the most influential I/O counters, which are used
% to retrieve relevant entries from a benchmark-grounded Knowledge Base.
% The LLM then generates specific, validated code modifications.
%
% [Results] TODO: Fill in after evaluation
%
% [Impact] TODO: Fill in after evaluation
